\begin{abstract}
The Indonesian clitic \id{-nya} is one of the highest-frequency morphological items in the language and serves at least four distinct functions (possessive, definite, anaphoric, and nominalising), yet standard Indonesian information-retrieval pipelines strip it as a generic suffix. We test whether and how this preprocessing default affects modern retrieval effectiveness through a controlled five-condition study on the Indonesian split of MIRACL: (i)~keep, (ii)~naive regex strip, (iii)~dictionary-aware Sastrawi clitic strip, (iv)~sentinel masking, and (v)~oracle anaphora resolution on a manually annotated subset. We evaluate both sparse (BM25) and dense (BGE-m3) retrievers and report nDCG@10 and MRR@100 across all conditions, with a paired Wilcoxon signed-rank protocol and a Bonferroni-corrected pairwise post-hoc. \todo{Insert headline result once experiments complete.} Our findings inform the preprocessing decisions made by every Indonesian retrieval system and bound the upper-effectiveness ceiling of future anaphora-aware preprocessing methods.
\end{abstract}

\noindent\textbf{Keywords:} information retrieval, Indonesian, morphology, clitic, anaphora, dense retrieval, MIRACL.
